\documentclass[10pt]{article}

% ======= Paquetes base =======
\usepackage[utf8]{inputenc}
\usepackage[T1]{fontenc}
\usepackage[spanish,es-nodecimaldot]{babel}
\usepackage{geometry}
\geometry{a5paper, top=1.5cm, bottom=1.5cm, left=2cm, right=1.5cm}
\usepackage{setspace}
\usepackage{titlesec}
\usepackage{enumitem}
\usepackage{graphicx}
\usepackage{float}
\usepackage{xcolor}
\usepackage{mathptmx}
\usepackage{caption}
\captionsetup{font=footnotesize}

% ======= Interlineado =======
\setstretch{1.0}
\setlength{\parindent}{1.25cm}
\setlength{\parskip}{0pt}

% ======= Títulos =======
\titleformat{\section}{\bfseries\fontsize{10}{12}\selectfont}{\thesection.}{0.6em}{}
\titleformat{\subsection}{\bfseries\fontsize{10}{12}\selectfont}{\thesubsection.}{0.6em}{}

\usepackage{hyperref}
\hypersetup{colorlinks=true, linkcolor=black, urlcolor=blue, citecolor=black}

\begin{document}

% =======================
% ENCABEZADO
% =======================
\begin{center}
  {\bfseries\fontsize{11}{13}\selectfont
  Sistema Omnicanal y Plataforma Progresiva para la Gestión\\
  de Mesas, Órdenes y Pagos en Restaurantes}\\[4pt]
  {\fontsize{9}{11}\selectfont
  Natalia Anaya Ramírez \quad Diego Villagrán Salazar}
\end{center}

\vspace{6pt}

% =======================
% METODOLOGÍA — SPRINTS
% =======================
\section{Metodología y Planificación}

El desarrollo sigue una metodología ágil basada en Scrum con sprints de duración fija (2--3 semanas cada uno), que permiten entregar funcionalidades incrementales y validarlas de manera continua. El proyecto se estructura en \textbf{12 sprints}:

\begin{enumerate}[leftmargin=*, itemsep=1pt, topsep=2pt]
  \item \textbf{Sprint 1:} Analítica de negocio y Arquitectura
  \item \textbf{Sprint 2:} Diseño de la base de datos
  \item \textbf{Sprint 3:} Desarrollo de la base de datos
  \item \textbf{Sprint 4:} Diseño de vistas y mockups
  \item \textbf{Sprint 5:} Diseño y arquitectura de backend
  \item \textbf{Sprint 6:} Definición de KPI y vistas de Dashboard
  \item \textbf{Sprint 7:} Implementación del dashboard
  \item \textbf{Sprint 8:} Desarrollo del PAW mesero
  \item \textbf{Sprint 9:} Desarrollo del PAW comensal
  \item \textbf{Sprint 10:} Desarrollo del PAW cocina
  \item \textbf{Sprint 11:} Integración de la API de pagos
  \item \textbf{Sprint 12:} Pruebas
\end{enumerate}

Se generará un tablero Kanban en GitHub con los estatus: \emph{inicializado}, \emph{en progreso}, \emph{en revisión} y \emph{finalizado}. Ninguno de estos estatus es final; puede volver a cualquier otro estado excepto inicializado. Pueden agregarse nuevas tareas dentro del sprint si es necesario.

% =======================
% CRONOGRAMA
% =======================
\section{Cronograma de Actividades}

Para la ejecución exitosa del proyecto se ha desarrollado un cronograma detallado que abarca el período de \textbf{febrero a diciembre de 2026}. La planificación considera las dependencias entre tareas, los recursos disponibles y los hitos críticos, asegurando una progresión ordenada desde la fase de análisis hasta la implementación completa del sistema omnicanal.

El cronograma está dividido en dos diagramas de Gantt que muestran la distribución de actividades y responsabilidades entre los miembros del equipo:

\begin{figure}[H]
\centering
\includegraphics[width=\linewidth]{Nat.png}
\caption{Cronograma de actividades — Natalia (Febrero–Diciembre 2026).}
\label{fig:gantt-nat}
\end{figure}

\begin{figure}[H]
\centering
\includegraphics[width=\linewidth]{Diego.png}
\caption{Cronograma de actividades — Diego (Febrero–Diciembre 2026).}
\label{fig:gantt-diego}
\end{figure}

\end{document}